\subsection{Verbindungs- und Ortsvektoren}
\label{subsec:analytic_geometry_connecting_and_position_vectors}

Mit einem Verbindungsvektor beschreiben wir die relative Lage zweier Punkte. 

\begin{DefinitionBox}{Verbindungsvektor}{analytic_geometry_connecting_vector}
\noindent
\begin{minipage}[c]{0.54\linewidth}
  Der \emph{Verbindungsvektor} \EnglishTerm{displacement vector} vom Punkt
  $P(p_1,p_2,p_3)$ zum Punkt $Q(q_1,q_2,q_3)$ ist
  \[
    \textcolor{THMBlue}{\overrightarrow{PQ}}
    =\begin{pmatrix}
       q_1-p_1\\ q_2-p_2\\ q_3-p_3
     \end{pmatrix}.
  \]
  Sein Aufpunkt ist $P$, seine Pfeilspitze liegt in $Q$.
\end{minipage}\hfill
\begin{minipage}[c]{0.42\linewidth}
  \centering
  \CachedTikzFigure{content/resources/chapters/analytic_geometry/tikz/connecting_position_vectors.tex}
\end{minipage}
\end{DefinitionBox}

\begin{KeyPointBox}{Abstand zweier Punkte}{analytic_geometry_point_distance}
  Der Abstand ergibt sich aus der Länge (Betrag) des Verbindungsvektors:
  \[
    \left|\textcolor{THMBlue}{\overrightarrow{PQ}}\right|
    =\left|\begin{pmatrix}
      q_1-p_1\\ q_2-p_2\\ q_3-p_3
    \end{pmatrix}\right|
    =\sqrt{
      (q_1-p_1)^2+(q_2-p_2)^2+(q_3-p_3)^2}.
  \]
\end{KeyPointBox}

\begin{ExampleBox}{Verbindungsvektor und Punktabstand}{analytic_geometry_point_distance_example}
\noindent
\begin{minipage}[c]{0.54\linewidth}
  Gegeben sind $P(1,1)$ und $Q(3,3)$. Der Verbindungsvektor lautet
  \[
    \overrightarrow{PQ}
    =\begin{ClozeContinuation}[11]{2}
      \begin{pmatrix}3-1\\3-1\end{pmatrix}
      =\begin{pmatrix}2\\2\end{pmatrix}
    \end{ClozeContinuation}.
  \]
  Bezeichnen wir den Abstand in diesem Beispiel mit $d$. Dann gilt
  \[
    d
    =\left|\overrightarrow{PQ}\right|
    =\Cloze{7cm}{\sqrt{2^2+2^2}=\sqrt 8=2\sqrt 2}.
  \]
  Die Rechnung bestimmt also zugleich Richtung und Länge der Verbindung.
\end{minipage}\hfill
\begin{minipage}[c]{0.40\linewidth}
  \centering
  \CachedTikzFigure{content/resources/chapters/analytic_geometry/tikz/point_distance_2d.tex}
\end{minipage}
\end{ExampleBox}

Zur Beschreibung der Lage eines Punktes $P$ im Raum betrachten wir den
Verbindungsvektor vom Ursprung $O(0,0,0)$ zu diesem Punkt $P$, 
den sogenannten Ortsvektor von $P$.
\begin{DefinitionBox}{Ortsvektor}{analytic_geometry_position_vector}
\noindent
\begin{minipage}[c]{0.54\linewidth}
  Der \emph{Ortsvektor} \EnglishTerm{position vector} eines Punktes
  $P(x_1,x_2,x_3)$ ist der Verbindungsvektor vom Ursprung
  $O(0,0,0)$ zu $P$:
  \[
    \overrightarrow{OP}=\begin{pmatrix}x_1\\x_2\\x_3\end{pmatrix}.
  \]
  Sein Aufpunkt ist immer der Ursprung.
\end{minipage}\hfill
\begin{minipage}[c]{0.42\linewidth}
  \centering
  \CachedTikzFigure{content/resources/chapters/analytic_geometry/tikz/position_vector_coordinates.tex}
\end{minipage}
\end{DefinitionBox}

\begin{ExampleBox}{Verschiebevektor und Ortsvektor}{analytic_geometry_displacement_and_position_vector_example}
  In diesem Beispiel zeichen wir den gleichen Vektor einmal als Verschiebevektor und einmal als Ortsvektor:
  \begin{itemize}[leftmargin=1.4em,itemsep=0.45em,topsep=0.35em]
    \item Der Verschiebevektor
      $\vec a=\begin{pmatrix}1\\2\end{pmatrix}$ kann an jedem Aufpunkt
      eingezeichnet werden.
    \item Der Ortsvektor
      $\overrightarrow{OP}=\begin{pmatrix}1\\2\end{pmatrix}$ beginnt
      dagegen zwingend im Ursprung.
  \end{itemize}

  \begin{center}
    \vspace{-0.45em}
    \centering
    \CachedTikzFigure{content/resources/chapters/analytic_geometry/tikz/displacement_and_position_vector.tex}
    \vspace{-0.55em}
  \end{center}
  %
  Jeder Verschiebevektor kann parallel verschoben werden, wogegen eine Ortsvektor immer vom Ursprung ausgeht.
\end{ExampleBox}
